\section{ساخت، \lr{Packaging} و \lr{Testing}}
\label{sec:book-manufacturing}
\subsection{\lr{Tapeout}، \lr{Mask} و تقسیم هزینه ساخت}
پس از نهایی‌شدن طراحی و بررسی‌های لازم، توصیف هندسه ماسک‌ها برای ساخت تحویل
داده می‌شود. این مرحله \lr{Tapeout} نام دارد. نام آن از روش تاریخی تحویل داده روی
نوار مغناطیسی آمده است. کتاب قالب‌های \lr{Caltech Interchange Format} یا
\lr{CIF} و \lr{Calma GDS II Stream Format} یا \lr{GDS} را معرفی می‌کند
\cite[بخش \secnum{۱.۱۲}]{weste}.

\lr{Mask} الگوی لازم برای انتقال هندسه یک مرحله از فرآیند را فراهم می‌کند.
کتاب روش ساخت الگوی کروم روی شیشه با پرتو الکترونی را توضیح می‌دهد.
هزینه ماسک و تجهیز خط تولید، بخش مهمی از هزینه ثابت ساخت است.
متن سال ۲۰۱۱ برای مثال‌های \qty{180}{\nano\metre} و \qty{65}{\nano\metre}
به‌ترتیب حدود ۲۵۰ هزار و ۳ میلیون دلار هزینه مجموعه ماسک ذکر می‌کند.
این مقادیر، مثال تاریخی همان کتاب‌اند.

در \lr{Multi-project Wafer}، چند طراحی کوچک از مجموعه مشترکی از ماسک‌ها استفاده
می‌کنند تا هزینه ثابت میان مشتریان تقسیم شود. کتاب خدمات \lr{MOSIS}،
\lr{EUROPRACTICE} و \lr{VDEC} را در زمینه زمان انتشار معرفی می‌کند.
در مدل آموزشی فرضی، اگر هزینه ثابت مشترک \(K\) میان \(N\) طراحی با سهم برابر
تقسیم شود، سهم هر طراحی \(K/N\) است. هزینه عملی به سهم سطح، فرآیند، تعداد نمونه،
\lr{Packaging} و شرایط خدمات وابسته است.

\lr{Fab} کارخانه ساخت و \lr{Foundry} تأمین‌کننده خدمات ساخت است.
شرکت \lr{Fabless} طراحی را انجام می‌دهد و ساخت را به کارخانه طرف قرارداد می‌سپارد.
کتاب هزینه بالای تأسیس کارخانه و اهمیت اقتصادی این تقسیم کار را توضیح می‌دهد.

\subsection{از \lr{Wafer} تا \lr{Die}}
روی یک \lr{Wafer} سیلیکونی، تعداد زیادی نسخه مدار هم‌زمان ساخته می‌شوند.
مراحل تکراری رسوب‌دهی، ماسک‌گذاری، \lr{Etch} و \lr{Implant} همان هندسه و
ساختار افزاره را ایجاد می‌کنند که در بخش ساخت مقدماتی آموختیم.
کتاب قطرهای \qty{150}{\milli\metre} تا \qty{300}{\milli\metre} و زمان بازگشت
تا حدود ۱۰ هفته را برای زمینه صنعتی زمان خود شرح می‌دهد \cite[صفحه ۵۴]{weste}.

در \lr{Clean Room}، فیلتر هوا و پوشش کارکنان، ورود ذرات آلاینده را کاهش می‌دهد.
یک ذره کوچک می‌تواند بخشی از الگو را خراب کند و اتصال باز یا اتصال کوتاه بسازد.
پس از فرآوری، \lr{Wafer} به قطعه‌های \lr{Die} بریده می‌شود. هندسه نهایی قطعه،
تعداد و محل اتصال‌ها و نیاز حرارتی، در انتخاب \lr{Package} نقش دارند.

\subsection{\lr{Wire Bonding} و \lr{Flip-chip}}
\lr{Package} اتصال قابل‌استفاده با محیط بیرونی و حفاظت مکانیکی قطعه را فراهم
می‌کند و در دفع گرما و رفتار الکتریکی نیز نقش دارد. مثال \lr{MIPS} کتاب، قطعه
\qty{1.5}{\milli\metre} در \qty{1.5}{\milli\metre} را در \lr{Dual-inline Package}
یا \lr{DIP} با ۴۰ پایه نشان می‌دهد. در \lr{Wire Bonding}، سیم‌های نازک طلا
\lr{Pad}های قطعه را به \lr{Lead Frame} بسته وصل می‌کنند
\cite[شکل \secnum{۱.۷۳} و صفحه ۵۵]{weste}.

در \lr{Flip-chip}، اتصال‌های لحیمی کوچک روی سطح قطعه ایجاد می‌شوند و قطعه
روبه‌پایین به زیرلایه متصل می‌شود. امکان توزیع اتصال‌ها در سطح تراشه، ظرفیت اتصال
را افزایش می‌دهد. حذف سیم \lr{Bond Wire} سهم القایی آن سیم را از مسیر برمی‌دارد.
رفتار الکتریکی نهایی همچنان به کل مسیر اتصال و بسته وابسته است.
کتاب معیارهای هزینه، تعداد پایه، پهنای باند پایه، انتقال توان و قابلیت اطمینان را
برای مقایسه بسته‌ها مطرح می‌کند و جزئیات را به فصل ۱۳ ارجاع می‌دهد.

\subsection{آزمون قطعه ساخته‌شده و \lr{BIST}}
\lr{Testing} قطعه ساخته‌شده را با معیار پذیرش می‌سنجد تا نمونه‌های دارای عیب
شناسایی شوند. حتی با طراحی صحیح، نقص ساخت یا آلودگی می‌تواند قطعه معیوب ایجاد کند.
آزمون به دسترسی ورودی‌ها، مشاهده پاسخ‌ها، سرعت موردنظر و هزینه زمانی دستگاه
آزمون وابسته است \cite[صفحه ۵۵]{weste}.

در \lr{Built-in Self-test} یا \lr{BIST}، مدارهای کمکی روی تراشه بخشی از تولید
تحریک و بررسی پاسخ را انجام می‌دهند. هدف معرفی‌شده در کتاب کاهش زمان لازم
دستگاه آزمون خارجی است. این انتخاب، منطق و مساحت کمکی می‌خواهد.
پوشش عیب به روش تولید الگو و تحلیل پاسخ وابسته است.
معیار آزمون باید با عیب‌هایی که قصد کشفشان را داریم متناسب باشد.

برای تحلیل اقتصادی ساده، \lr{Yield} را نسبت تعداد قطعات سالم به تعداد قطعات
ساخته‌شده تعریف می‌کنیم. اگر به فرض ۱۰۰۰ قطعه ساخته شود و ۸۵۰ قطعه معیار پذیرش
را برآورده کنند، \(Y=850/1000=0.85\)، یعنی ۸۵ درصد است.
اگر هزینه مفروض همین مرحله \(K\) باشد، هزینه سرشکن‌شده هر قطعه سالم در این مدل
\begin{equation}
 C_{\mathrm{good}}=\frac{K}{N_{\mathrm{good}}}
 =\frac{K}{YN_{\mathrm{total}}}
 \label{eq:impl-yield-cost}
\end{equation}
خواهد بود. \(N_{\mathrm{total}}\) تعداد کل و \(N_{\mathrm{good}}\) تعداد سالم است.
با ثابت‌بودن هزینه و تعداد ساخت، کاهش \lr{Yield} هزینه هر قطعه سالم را افزایش می‌دهد.
این مثال افزوده، مدل سرشکن‌کردن یک هزینه مشخص است. هزینه‌های مراحل دیگر باید
در مدل اقتصادی کامل خود لحاظ شوند.

\subsection{تمرین‌های پاسخ‌دار ساخت و آزمون}
\textbf{تمرین ۱.} هزینه فرضی یک مرحله برای ۱۰۰۰ قطعه، ۲۰ هزار واحد پول است.
با \lr{Yield} برابر ۸۰ درصد، هزینه همان مرحله برای هر قطعه سالم چقدر می‌شود؟
تعداد سالم \(0.8(1000)=800\) و هزینه واحد \(20000/800=25\) واحد پول است.
با \lr{Yield} برابر ۵۰ درصد، هزینه \(20000/500=40\) واحد پول می‌شود.

\textbf{تمرین ۲.} برای عیب خروجی گیرکرده روی صفر در \lr{Carry}، چه ورودی‌ای مناسب است؟
مثلاً \(a=b=1,c=0\) پاسخ مرجع یک دارد و عیب را در مدل گیرکردن خروجی آشکار می‌کند.
ورودی \(0,0,0\) پاسخ مرجع صفر دارد و این عیب خاص را آشکار نمی‌کند.
همان منطق انتخاب تحریک در \lr{Testbench} نیز کاربرد دارد.

\textbf{تمرین ۳.} اگر \lr{LVS} و آزمون عملکرد پیش از ساخت موفق باشند، قطعه ساخته‌شده
چه بررسی دیگری می‌خواهد؟ شرایط ساخت هر نمونه می‌تواند نقص ایجاد کند، بنابراین
آزمون قطعه با معیارهای عملکرد و کیفیت تعیین‌شده لازم است. صحت اتصال استخراج‌شده،
صحت رفتار مدل و سلامت نمونه فیزیکی، هر یک به شواهد مناسب خود وابسته‌اند.

\subsection{پیوند این بخش با مدل‌های فصل ۲}
برای محاسبه دقیق‌تر تأخیر باید بدانیم جریان ترانزیستور چگونه به
\(V_{GS}\)، \(V_{DS}\)، \(V_{BS}\)، \(W\) و \(L\) وابسته است.
برای محاسبه بار باید خازن‌های \lr{Gate}، \lr{Diffusion} و سیم را بشناسیم.
برای ارزیابی حالت منطقی پایدار نیز رابطه ورودی و خروجی \lr{Inverter} اهمیت دارد.
بخش‌های آغازین فصل ۲ این مدل‌های افزاره را می‌سازند و بخش‌های بعدی اثرهای غیرایده‌آل
و مشخصه‌های \lr{DC} را به آن‌ها اضافه می‌کنند.
\clearpage
